%----------------------------------------------------------------------------------------
% GONIOMETRISCHE FORMULES. Alle opgaven zijn eigen werk.
%
% Opbouw: Start, leerstof (theorie, werkwijze of uitleg, voorbeeld, korte oefening), en op een
% nieuwe pagina de Opdrachten (door elkaar, stijgende moeilijkheid), Koppel en Test jezelf.
% Labels van dit hoofdstuk beginnen met h2-.
%----------------------------------------------------------------------------------------

\mychapter{2}{Goniometrische formules}

\lkr{Leerplan B+S: LPD 13-16 (goniometrische functies), het deel over de som- en verschilformules en de verdubbelingsformules. Het bewijs van de verschilformule voor de cosinus is een mooi voorbeeld van een bewijs met een afstand: laat leerlingen bij elke regel zeggen waarom ze klopt. Sectie 3 (sinus en cosinus samenvoegen) bereidt de algemene sinusfunctie in hoofdstuk 3 voor.}

Is $\sin(\alpha + \beta)$ gelijk aan $\sin\alpha + \sin\beta$? Probeer $\alpha = \beta = \frac{\pi}{2}$: links staat $\sin\pi = 0$, rechts $1 + 1 = 2$. Nee dus. In dit hoofdstuk zoeken we wat $\sin(\alpha + \beta)$ en $\cos(\alpha + \beta)$ wel zijn. Daarmee bereken je exacte waarden voor nieuwe hoeken, en vereenvoudig je uitdrukkingen.

Alles volgt uit één formule: de cosinus van een verschil. Die bewijzen we met een afstand op de eenheidscirkel. Daarna leiden we de andere formules er één voor één uit af, telkens in een paar regels.

\begin{oefeningen}[Start]
\oef{h2-start-1}

\oef{h2-verken}
\end{oefeningen}

%========================================================================================
\section{Som- en verschilformules}

\subsection{De cosinus van een verschil}

Neem op de eenheidscirkel het punt $P$ bij de hoek $\alpha$ en het punt $Q$ bij de hoek $\beta$. De hoek tussen $OP$ en $OQ$ is $\alpha - \beta$. Draai je beide punten samen rond $O$, dan verandert die hoek niet, en de afstand $|PQ|$ ook niet. We berekenen die afstand op twee manieren: met de coördinaten van $P$ en $Q$, en na een draaiing die $Q$ naar $(1, 0)$ brengt.

\begin{theorie}
\begin{stelling}{Verschilformule voor de cosinus}
Voor alle hoeken $\alpha$ en $\beta$:

\vspace{1mm}
\qquad $\formule{\cos(\alpha - \beta) = \cos\alpha\cos\beta + \sin\alpha\sin\beta}$
\end{stelling}
\begin{bewijs}
\naast{%
Neem $P(\cos\alpha, \sin\alpha)$ en $Q(\cos\beta, \sin\beta)$. Met de afstandsformule is
\bewijsregel{|PQ|^2 = (\cos\alpha - \cos\beta)^2 + (\sin\alpha - \sin\beta)^2 .}
Werk de kwadraten uit en gebruik twee keer de grondformule:
\bewijsregel{|PQ|^2 = 2 - 2(\cos\alpha\cos\beta + \sin\alpha\sin\beta) .}
Draai beide punten over $-\beta$: $Q$ gaat naar $Q'(1, 0)$ en $P$ naar $P'\bigl(\cos(\alpha - \beta), \sin(\alpha - \beta)\bigr)$. Een draaiing bewaart afstanden, dus $|P'Q'| = |PQ|$, en
\bewijsregel{|P'Q'|^2 = \bigl(\cos(\alpha - \beta) - 1\bigr)^2 + \sin^2(\alpha - \beta) = 2 - 2\cos(\alpha - \beta) .}}{%
\begin{tikzpicture}[scale=1.3,font=\scriptsize]
  \eenheidscirkel
  \draw[accent] (0,0) -- (115:1) (0,0) -- (40:1);
  \draw[accent,thick] (115:1) -- (40:1);
  \fill[accent] (115:1) circle (.035) node[above left] {$P$};
  \fill[accent] (40:1) circle (.035) node[right] {$Q$};
  \draw[accent,dashed] (0,0) -- (75:1);
  \draw[accent,dashed,thick] (75:1) -- (1,0);
  \fill[accent] (75:1) circle (.035) node[above] {$P'$};
  \fill[accent] (1,0) circle (.035) node[below right] {$Q'$};
\end{tikzpicture}}
De twee uitdrukkingen voor dezelfde afstand zijn gelijk:
\bewijsregel{2 - 2\cos(\alpha - \beta) = 2 - 2(\cos\alpha\cos\beta + \sin\alpha\sin\beta) ,}
waaruit de formule volgt.
\end{bewijs}
\end{theorie}

De formule zegt dat je de cosinus van een verschil kunt berekenen uit sinus en cosinus van de twee hoeken apart. Let op: ze is niet $\cos\alpha - \cos\beta$. Neem $\alpha = \beta$: links staat $\cos 0 = 1$, terwijl $\cos\alpha - \cos\alpha = 0$.

\voorbeeld{}{Bepaal exact $\cos\frac{\pi}{12}$.\label{vb:h2-exact}}{%
\uitleg{$\frac{\pi}{12}$ is geen bijzondere hoek, maar wel het verschil van twee: $\frac{\pi}{3} - \frac{\pi}{4} = \frac{4\pi - 3\pi}{12} = \frac{\pi}{12}$. Van die twee hoeken kennen we sinus en cosinus.}
\begin{stappenbreed}
\cos\tfrac{\pi}{12} = \cos\tfrac{\pi}{3}\cos\tfrac{\pi}{4} + \sin\tfrac{\pi}{3}\sin\tfrac{\pi}{4} & Verschilformule met $\alpha = \frac{\pi}{3}$ en $\beta = \frac{\pi}{4}$. \\[1ex]
= \tfrac12 \cdot \tfrac{\sqrt2}{2} + \tfrac{\sqrt3}{2} \cdot \tfrac{\sqrt2}{2} & Bijzondere waarden. \\[1ex]
= \tfrac{\sqrt2 + \sqrt6}{4} & Controle: $\approx 0{,}9659 \approx \cos 15^\circ$.
\end{stappenbreed}}

\begin{werkwijze}
Schrijf de hoek als som of verschil van twee hoeken uit de tabel. Bij een noemer $12$ lukt dat bijvoorbeeld met $\frac{\pi}{3}$ en $\frac{\pi}{4}$, of $\frac{\pi}{4}$ en $\frac{\pi}{6}$. Kies de passende formule, vul de bijzondere waarden in en vereenvoudig. Controleer met het rekentoestel.
\end{werkwijze}

\begin{oefeningen}
\oef{h2-cosverschil-k}
\end{oefeningen}

\needspace{14\baselineskip}
\subsection{De cosinus van een som}

Een som is een verschil met een tegengestelde hoek: $\alpha + \beta = \alpha - (-\beta)$. Zo volgt de somformule uit de verschilformule, met de tegengestelde hoeken uit hoofdstuk 1.

\begin{theorie}
\begin{stelling}{Somformule voor de cosinus}
Voor alle hoeken $\alpha$ en $\beta$:

\vspace{1mm}
\qquad $\formule{\cos(\alpha + \beta) = \cos\alpha\cos\beta - \sin\alpha\sin\beta}$
\end{stelling}
\begin{bewijs}
Pas de verschilformule toe op $\alpha$ en $-\beta$, en gebruik $\cos(-\beta) = \cos\beta$ en $\sin(-\beta) = -\sin\beta$:
\bewijsregel{\cos(\alpha + \beta) = \cos\bigl(\alpha - (-\beta)\bigr) = \cos\alpha\cos(-\beta) + \sin\alpha\sin(-\beta) = \cos\alpha\cos\beta - \sin\alpha\sin\beta .}
\end{bewijs}
\end{theorie}

Som en verschil verschillen alleen in het teken in het midden, en dat teken is net omgekeerd: een som in de hoek geeft een min, een verschil een plus.

\voorbeeld{}{Bepaal exact $\cos\frac{7\pi}{12}$.\label{vb:h2-cossom}}{%
\uitleg{$\frac{7\pi}{12} = \frac{4\pi}{12} + \frac{3\pi}{12} = \frac{\pi}{3} + \frac{\pi}{4}$: een som van twee hoeken uit de tabel.}
\begin{stappenbreed}
\cos\tfrac{7\pi}{12} = \cos\tfrac{\pi}{3}\cos\tfrac{\pi}{4} - \sin\tfrac{\pi}{3}\sin\tfrac{\pi}{4} & Somformule. \\[1ex]
= \tfrac12 \cdot \tfrac{\sqrt2}{2} - \tfrac{\sqrt3}{2} \cdot \tfrac{\sqrt2}{2} = \tfrac{\sqrt2 - \sqrt6}{4} & Bijzondere waarden.
\end{stappenbreed}
Controle: de uitkomst is negatief, en dat klopt, want $\frac{7\pi}{12}$ ligt in het tweede kwadrant. Het rekentoestel geeft $\cos 105^\circ \approx -0{,}2588$.}

\begin{oefeningen}
\oef{h2-cossom-k}
\end{oefeningen}

\needspace{16\baselineskip}
\subsection{De sinus van een som}

Voor de sinus gebruiken we de complementaire hoeken: $\sin\gamma = \cos\bigl(\frac{\pi}{2} - \gamma\bigr)$. Zo wordt de sinus van een som de cosinus van een verschil, en die formule kennen we al.

\begin{theorie}
\begin{stelling}{Somformule voor de sinus}
Voor alle hoeken $\alpha$ en $\beta$:

\vspace{1mm}
\qquad $\formule{\sin(\alpha + \beta) = \sin\alpha\cos\beta + \cos\alpha\sin\beta}$
\end{stelling}
\begin{bewijs}
Schrijf de sinus als cosinus van de complementaire hoek, en groepeer:
\bewijsregel{\sin(\alpha + \beta) = \cos\bigl(\tfrac{\pi}{2} - \alpha - \beta\bigr) = \cos\bigl((\tfrac{\pi}{2} - \alpha) - \beta\bigr) .}
De verschilformule voor de cosinus geeft
\bewijsregel{\sin(\alpha + \beta) = \cos\bigl(\tfrac{\pi}{2} - \alpha\bigr)\cos\beta + \sin\bigl(\tfrac{\pi}{2} - \alpha\bigr)\sin\beta = \sin\alpha\cos\beta + \cos\alpha\sin\beta ,}
weer met complementaire hoeken.
\end{bewijs}
\end{theorie}

Bij de sinus wisselen sinus en cosinus elkaar af in de twee termen, en het teken blijft hetzelfde als in de hoek.

\voorbeeld{}{Gegeven: $\sin\alpha = \frac35$ en $\cos\beta = \frac{5}{13}$, met $\alpha$ en $\beta$ in het eerste kwadrant. Bereken \mbox{$\sin(\alpha + \beta)$}.\label{vb:h2-som}}{%
\uitleg{De somformule heeft van elke hoek de sinus én de cosinus nodig. We zoeken eerst de ontbrekende cosinus en sinus met de grondformule. In het eerste kwadrant zijn ze positief.}
\begin{stappenbreed}
\cos\alpha = \sqrt{1 - \tfrac{9}{25}} = \tfrac45 & Grondformule, positieve wortel. \\[1ex]
\sin\beta = \sqrt{1 - \tfrac{25}{169}} = \tfrac{12}{13} & Idem. \\[1ex]
\sin(\alpha + \beta) = \tfrac35 \cdot \tfrac{5}{13} + \tfrac45 \cdot \tfrac{12}{13} = \tfrac{15 + 48}{65} = \tfrac{63}{65} & Somformule.
\end{stappenbreed}}

\begin{oefeningen}
\oef{h2-sinsom-k}
\end{oefeningen}

\needspace{14\baselineskip}
\subsection{De sinus van een verschil}

Net zoals bij de cosinus vervang je $\beta$ door $-\beta$.

\begin{theorie}
\begin{stelling}{Verschilformule voor de sinus}
Voor alle hoeken $\alpha$ en $\beta$:

\vspace{1mm}
\qquad $\formule{\sin(\alpha - \beta) = \sin\alpha\cos\beta - \cos\alpha\sin\beta}$
\end{stelling}
\begin{bewijs}
Pas de somformule toe op $\alpha$ en $-\beta$:
\bewijsregel{\sin(\alpha - \beta) = \sin\alpha\cos(-\beta) + \cos\alpha\sin(-\beta) = \sin\alpha\cos\beta - \cos\alpha\sin\beta .}
\end{bewijs}
\end{theorie}

\voorbeeld{}{Bepaal exact $\sin\frac{\pi}{12}$.\label{vb:h2-sinverschil}}{%
\uitleg{Zoals in voorbeeld~\ref{vb:h2-exact} is $\frac{\pi}{12} = \frac{\pi}{3} - \frac{\pi}{4}$. Nu gaat het om de sinus, dus we gebruiken de verschilformule voor de sinus.}
\begin{stappenbreed}
\sin\tfrac{\pi}{12} = \sin\tfrac{\pi}{3}\cos\tfrac{\pi}{4} - \cos\tfrac{\pi}{3}\sin\tfrac{\pi}{4} & Verschilformule. \\[1ex]
= \tfrac{\sqrt3}{2} \cdot \tfrac{\sqrt2}{2} - \tfrac12 \cdot \tfrac{\sqrt2}{2} = \tfrac{\sqrt6 - \sqrt2}{4} & Controle: $\approx 0{,}2588 \approx \sin 15^\circ$.
\end{stappenbreed}}

\begin{oefeningen}
\oef{h2-sinverschil-k}
\end{oefeningen}

\needspace{16\baselineskip}
\subsection{De tangens van een som en een verschil}

De tangens is het quotiënt van sinus en cosinus. Deel je de somformule voor de sinus door die voor de cosinus, dan krijg je een formule met alleen tangensen. Ze geldt zolang alle tangensen in de formule bestaan.

\begin{theorie}
\begin{stelling}{Somformule voor de tangens}
\qquad $\formule{\tan(\alpha + \beta) = \frac{\tan\alpha + \tan\beta}{1 - \tan\alpha\tan\beta}}$

\vspace{1mm}
als $\tan\alpha$, $\tan\beta$ en $\tan(\alpha + \beta)$ bestaan.
\end{stelling}
\begin{bewijs}
Volgens de somformules is
\bewijsregel{\tan(\alpha + \beta) = \frac{\sin(\alpha + \beta)}{\cos(\alpha + \beta)} = \frac{\sin\alpha\cos\beta + \cos\alpha\sin\beta}{\cos\alpha\cos\beta - \sin\alpha\sin\beta} .}
Omdat $\tan\alpha$ en $\tan\beta$ bestaan, is $\cos\alpha\cos\beta \neq 0$. Deel teller en noemer door $\cos\alpha\cos\beta$:
\bewijsregel{\tan(\alpha + \beta) = \frac{\frac{\sin\alpha}{\cos\alpha} + \frac{\sin\beta}{\cos\beta}}{1 - \frac{\sin\alpha}{\cos\alpha} \cdot \frac{\sin\beta}{\cos\beta}} = \frac{\tan\alpha + \tan\beta}{1 - \tan\alpha\tan\beta} .}
\end{bewijs}
\begin{stelling}{Verschilformule voor de tangens}
\qquad $\formule{\tan(\alpha - \beta) = \frac{\tan\alpha - \tan\beta}{1 + \tan\alpha\tan\beta}}$

\vspace{1mm}
als $\tan\alpha$, $\tan\beta$ en $\tan(\alpha - \beta)$ bestaan.
\end{stelling}
\begin{bewijs}
Pas de somformule toe op $\alpha$ en $-\beta$, met $\tan(-\beta) = -\tan\beta$.
\end{bewijs}
\end{theorie}

Ken je van twee hoeken alleen de tangens, dan heb je zo de tangens van hun som, zonder eerst sinus en cosinus te zoeken.

\voorbeeld{}{Twee scherpe hoeken $\alpha$ en $\beta$ hebben $\tan\alpha = \frac15$ en $\tan\beta = \frac23$. Bereken $\alpha + \beta$.\label{vb:h2-tansom}}{%
\uitleg{We kennen alleen de tangensen, dus de somformule voor de tangens ligt voor de hand. Daarna zoeken we de hoek met die tangens.}
\begin{stappenbreed}
\tan(\alpha + \beta) = \dfrac{\frac15 + \frac23}{1 - \frac15 \cdot \frac23} = \dfrac{\frac{13}{15}}{\frac{13}{15}} = 1 & Somformule voor de tangens. \\[3ex]
0 < \alpha + \beta < \pi & Twee scherpe hoeken samen. \\[1ex]
\alpha + \beta = \tfrac{\pi}{4} & De enige hoek in $]0, \pi[$ met tangens $1$.
\end{stappenbreed}}

\begin{oefeningen}
\oef{h2-tan-k}
\end{oefeningen}

\subsection{Welke formule?}

Bij een opgave moet je zelf kiezen welke formule past. Kijk naar wat gevraagd wordt (sinus, cosinus of tangens) en naar de hoek (som of verschil). Ontbreekt een sinus of cosinus, zoek die dan eerst met de grondformule, en let op het kwadrant.

\begin{oefeningen}
\oef{h2-exactk}

\oef{h2-som}
\end{oefeningen}

%========================================================================================
\section{Verdubbelingsformules}

Neem $\beta = \alpha$ in de somformules. Dan krijg je formules voor de dubbele hoek $2\alpha$. Ook hier geldt: de sinus van het dubbele is niet het dubbele van de sinus. Voor $\alpha = \frac{\pi}{2}$ is $\sin 2\alpha = \sin\pi = 0$, terwijl $2\sin\frac{\pi}{2} = 2$.

\subsection{De sinus van de dubbele hoek}

\begin{theorie}
\begin{stelling}{Verdubbelingsformule voor de sinus}
Voor elke hoek $\alpha$: \quad $\formule{\sin 2\alpha = 2\sin\alpha\cos\alpha}$
\end{stelling}
\begin{bewijs}
Neem $\beta = \alpha$ in de somformule voor de sinus:
\bewijsregel{\sin 2\alpha = \sin(\alpha + \alpha) = \sin\alpha\cos\alpha + \cos\alpha\sin\alpha = 2\sin\alpha\cos\alpha .}
\end{bewijs}
\end{theorie}

Lees de formule ook van rechts naar links: een product $\sin\alpha\cos\alpha$ is de helft van $\sin 2\alpha$.

\voorbeeld{}{Bereken exact $\sin\frac{\pi}{12}\cos\frac{\pi}{12}$.\label{vb:h2-sindubbel}}{%
\uitleg{Sinus en cosinus van dezelfde hoek, vermenigvuldigd: dat is de helft van een verdubbelingsformule. Zo hoeven we $\sin\frac{\pi}{12}$ en $\cos\frac{\pi}{12}$ niet apart te kennen.}
\begin{stappenbreed}
\sin\tfrac{\pi}{12}\cos\tfrac{\pi}{12} = \tfrac12 \cdot 2\sin\tfrac{\pi}{12}\cos\tfrac{\pi}{12} & Vermenigvuldig met $2$ en deel door $2$. \\[1ex]
= \tfrac12 \sin\tfrac{2\pi}{12} = \tfrac12 \sin\tfrac{\pi}{6} = \tfrac14 & Verdubbelingsformule, met $\alpha = \frac{\pi}{12}$.
\end{stappenbreed}
Controle met de waarden uit de vorige voorbeelden: $\frac{\sqrt6 - \sqrt2}{4} \cdot \frac{\sqrt6 + \sqrt2}{4} = \frac{6 - 2}{16} = \frac14$.}

\begin{oefeningen}
\oef{h2-sindubbel-k}
\end{oefeningen}

\needspace{14\baselineskip}
\subsection{De cosinus van de dubbele hoek}

Voor de cosinus geeft de somformule een eerste vorm. Met de grondformule kun je daarin $\sin^2\alpha$ of $\cos^2\alpha$ vervangen. Zo krijg je drie vormen, en je kiest telkens de vorm die het best past.

\begin{theorie}
\begin{stelling}{Verdubbelingsformule voor de cosinus}
Voor elke hoek $\alpha$:

\vspace{1mm}
\qquad $\formule{\cos 2\alpha = \cos^2\alpha - \sin^2\alpha = 2\cos^2\alpha - 1 = 1 - 2\sin^2\alpha}$
\end{stelling}
\begin{bewijs}
Neem $\beta = \alpha$ in de somformule voor de cosinus:
\bewijsregel{\cos 2\alpha = \cos\alpha\cos\alpha - \sin\alpha\sin\alpha = \cos^2\alpha - \sin^2\alpha .}
Vervang $\sin^2\alpha$ door $1 - \cos^2\alpha$, dan is $\cos 2\alpha = 2\cos^2\alpha - 1$. Vervang je $\cos^2\alpha$ door $1 - \sin^2\alpha$, dan is $\cos 2\alpha = 1 - 2\sin^2\alpha$.
\end{bewijs}
\end{theorie}

Ken je alleen $\sin\alpha$, dan is de laatste vorm het snelst. Ken je alleen $\cos\alpha$, dan de middelste.

\voorbeeld{}{Gegeven: $\sin\alpha = \frac35$ en $\frac{\pi}{2} < \alpha < \pi$. Bereken $\sin 2\alpha$ en $\cos 2\alpha$.\label{vb:h2-dubbel}}{%
\uitleg{Voor $\sin 2\alpha$ hebben we ook $\cos\alpha$ nodig: in het tweede kwadrant is $\cos\alpha = -\frac45$. Voor $\cos 2\alpha$ is de vorm met alleen $\sin\alpha$ het snelst.}
\begin{stappenbreed}
\sin 2\alpha = 2 \cdot \tfrac35 \cdot \bigl(-\tfrac45\bigr) = -\tfrac{24}{25} & Verdubbelingsformule voor de sinus. \\[1ex]
\cos 2\alpha = 1 - 2 \cdot \tfrac{9}{25} = \tfrac{7}{25} & Vorm met $\sin^2\alpha$. Controle: $\bigl(\tfrac{24}{25}\bigr)^2 + \bigl(\tfrac{7}{25}\bigr)^2 = 1$.
\end{stappenbreed}}

\begin{oefeningen}
\oef{h2-verdubbel}
\end{oefeningen}

\needspace{14\baselineskip}
\subsection{De tangens van de dubbele hoek}

\begin{theorie}
\begin{stelling}{Verdubbelingsformule voor de tangens}
\qquad $\formule{\tan 2\alpha = \frac{2\tan\alpha}{1 - \tan^2\alpha}}$ \quad als $\tan\alpha$ en $\tan 2\alpha$ bestaan.
\end{stelling}
\begin{bewijs}
Neem $\beta = \alpha$ in de somformule voor de tangens:
\bewijsregel{\tan 2\alpha = \frac{\tan\alpha + \tan\alpha}{1 - \tan\alpha \cdot \tan\alpha} = \frac{2\tan\alpha}{1 - \tan^2\alpha} .}
\end{bewijs}
\end{theorie}

\voorbeeld{}{Gegeven: $\tan\alpha = -2$ en $\frac{\pi}{2} < \alpha < \pi$. Bereken $\tan 2\alpha$.\label{vb:h2-tandubbel}}{%
\uitleg{Alleen de tangens is gegeven, en de tangens is gevraagd. De formule voor $\tan 2\alpha$ gebruikt alleen $\tan\alpha$.}
\begin{stappenbreed}
\tan 2\alpha = \dfrac{2 \cdot (-2)}{1 - (-2)^2} = \dfrac{-4}{-3} = \dfrac43 & Verdubbelingsformule voor de tangens.
\end{stappenbreed}
Controle met het kwadrant: $\pi < 2\alpha < 2\pi$. Een positieve tangens past bij het derde kwadrant, en $2\alpha$ ligt dus tussen $\pi$ en $\frac{3\pi}{2}$.}

\begin{oefeningen}
\oef{h2-dubbel}
\end{oefeningen}

\needspace{20\baselineskip}
\subsection{De formules van Carnot}

De vormen $2\cos^2\alpha - 1$ en $1 - 2\sin^2\alpha$ van $\cos 2\alpha$ kun je oplossen naar $\cos^2\alpha$ en $\sin^2\alpha$. Zo ruil je een kwadraat in voor een dubbele hoek. Dat heb je nodig om kwadraten weg te werken, bijvoorbeeld bij het integreren, en om sinus en cosinus van een halve hoek te berekenen.

\begin{theorie}
\begin{stelling}{Formule van Carnot voor de cosinus}
Voor elke hoek $\alpha$: \quad $\formule{\cos^2\alpha = \frac{1 + \cos 2\alpha}{2}}$
\end{stelling}
\begin{bewijs}
Los $\cos 2\alpha = 2\cos^2\alpha - 1$ op naar $\cos^2\alpha$: \quad $2\cos^2\alpha = 1 + \cos 2\alpha$.
\end{bewijs}
\begin{stelling}{Formule van Carnot voor de sinus}
Voor elke hoek $\alpha$: \quad $\formule{\sin^2\alpha = \frac{1 - \cos 2\alpha}{2}}$
\end{stelling}
\begin{bewijs}
Los $\cos 2\alpha = 1 - 2\sin^2\alpha$ op naar $\sin^2\alpha$: \quad $2\sin^2\alpha = 1 - \cos 2\alpha$.
\end{bewijs}
\end{theorie}

Met $\alpha = \frac{\theta}{2}$ geven de formules het kwadraat van sinus en cosinus van de halve hoek $\frac{\theta}{2}$. Het teken van de wortel kies je daarna met het kwadrant.

\voorbeeld{}{Bepaal exact $\cos\frac{5\pi}{8}$.\label{vb:h2-carnot}}{%
\uitleg{$\frac{5\pi}{8}$ is de helft van $\frac{5\pi}{4}$, en van die hoek kennen we de cosinus. Carnot geeft dus het kwadraat, het kwadrant geeft het teken.}
\begin{stappenbreed}
\cos^2\tfrac{5\pi}{8} = \dfrac{1 + \cos\frac{5\pi}{4}}{2} = \dfrac{1 - \frac{\sqrt2}{2}}{2} = \dfrac{2 - \sqrt2}{4} & Carnot voor de cosinus, met $\cos\frac{5\pi}{4} = -\frac{\sqrt2}{2}$. \\[3ex]
\cos\tfrac{5\pi}{8} = -\dfrac{\sqrt{2 - \sqrt2}}{2} & $\frac{5\pi}{8}$ ligt in het tweede kwadrant: de cosinus is negatief.
\end{stappenbreed}
Controle: $-\frac{\sqrt{2 - \sqrt2}}{2} \approx -0{,}3827 \approx \cos 112{,}5^\circ$.}

\voorbeeld{}{Bepaal exact $\sin\frac{\pi}{12}$ met een formule van Carnot, en vergelijk met voorbeeld~\ref{vb:h2-sinverschil}.\label{vb:h2-carnotsin}}{%
\uitleg{$\frac{\pi}{12}$ is de helft van $\frac{\pi}{6}$. Carnot voor de sinus geeft het kwadraat. Het resultaat ziet er anders uit dan in voorbeeld~\ref{vb:h2-sinverschil}: we gaan na dat het hetzelfde getal is.}
\begin{stappenbreed}
\sin^2\tfrac{\pi}{12} = \dfrac{1 - \cos\frac{\pi}{6}}{2} = \dfrac{2 - \sqrt3}{4} & Carnot voor de sinus. \\[3ex]
\sin\tfrac{\pi}{12} = \dfrac{\sqrt{2 - \sqrt3}}{2} & Eerste kwadrant: positief. \\[3ex]
\Bigl(\dfrac{\sqrt6 - \sqrt2}{4}\Bigr)^2 = \dfrac{8 - 2\sqrt{12}}{16} = \dfrac{2 - \sqrt3}{4} & Het kwadraat van de vorige uitkomst is hetzelfde. Beide zijn positief, dus ze zijn gelijk.
\end{stappenbreed}}

\begin{oefeningen}
\oef{h2-carnotk}
\end{oefeningen}

\subsection{Identiteiten bewijzen}

Met de verdubbelingsformules kun je uitdrukkingen vereenvoudigen en identiteiten bewijzen. De kunst is de juiste vorm te kiezen.

\voorbeeld{}{Bewijs: $\dfrac{1 - \cos 2x}{\sin 2x} = \tan x$ (voor $\sin 2x \neq 0$).\label{vb:h2-identiteit}}{%
\uitleg{Links staat de dubbele hoek, rechts de enkele. We schrijven dus alles met $x$. Voor $\cos 2x$ kiezen we de vorm die de $1$ wegwerkt.}
\begin{stappenbreed}
\frac{1 - \cos 2x}{\sin 2x} = \frac{1 - (1 - 2\sin^2 x)}{2\sin x\cos x} & Verdubbelingsformules. \\[2.5ex]
= \frac{2\sin^2 x}{2\sin x\cos x} = \frac{\sin x}{\cos x} = \tan x & Vereenvoudigen door $2\sin x$; dat is niet $0$, want $\sin 2x \neq 0$.
\end{stappenbreed}}

\begin{werkwijze}
Bij een identiteit: vertrek van één lid (meestal het ingewikkeldste) en werk het uit tot het andere. Schrijf alles met dezelfde hoek ($2x$ als $x$), en alles met $\sin$ en $\cos$. Kies de vorm van $\cos 2x$ die het meest vereenvoudigt.
\end{werkwijze}

\begin{oefeningen}
\oef{h2-identk}
\end{oefeningen}

%========================================================================================
\section{Sinus en cosinus samenvoegen}

Tel je een sinus en een cosinus op, dan krijg je opnieuw een golf met dezelfde periode, alleen hoger en verschoven. In de figuur is de som van $\sin x$ en $\sqrt3\cos x$ (streepjes) de volle kromme: een sinusgrafiek met hoogte $2$, over $\frac{\pi}{3}$ naar links geschoven.

\figuren{%
\begin{tikzpicture}[x=.8cm,y=.7cm,font=\scriptsize]
  \draw[black!14,very thin] (-1.0472,-2) grid[xstep=1.0472,ystep=1] (6.2832,2);
  \draw[->,grijs] (-1.4,0) -- (6.8,0) node[right,text=black] {$x$};
  \draw[->,grijs] (0,-2.4) -- (0,2.6) node[above,text=black] {$y$};
  \foreach \t/\n in {-1.0472/-\frac{\pi}{3},1.0472/\frac{\pi}{3},3.1416/\pi,5.236/\frac{5\pi}{3}}{\node[below,inner sep=1.5pt] at (\t,0) {$\n$};}
  \foreach \t in {-2,-1,1,2}{\node[left,inner sep=1.5pt] at (0,\t) {$\t$};}
  \draw[accent,dashed] plot[domain=-1.0472:6.2832,samples=100,smooth] (\x,{sin(\x r)});
  \draw[accent,dashed] plot[domain=-1.0472:6.2832,samples=100,smooth] (\x,{sqrt(3)*cos(\x r)});
  \draw[accent,very thick] plot[domain=-1.0472:6.2832,samples=100,smooth] (\x,{2*sin((\x+1.0472) r)});
  \node[accent,right] at (1.2,1.9) {$2\sin\bigl(x + \frac{\pi}{3}\bigr)$};
\end{tikzpicture}}

Dat is geen toeval: met de somformule voor de sinus, van rechts naar links gelezen, schrijf je elke som $a\sin x + b\cos x$ als één verschoven sinus.

\begin{theorie}
\begin{stelling}{Samenvoegen}
Voor $(a, b) \neq (0, 0)$ is \quad $\formule{a\sin x + b\cos x = R\sin(x + \varphi)}$

\vspace{1mm}
met \quad $\formule{R = \sqrt{a^2 + b^2}}$, \quad $\formule{\cos\varphi = \frac{a}{R}}$ \ en \ $\formule{\sin\varphi = \frac{b}{R}}$.
\end{stelling}
\begin{bewijs}
Omdat $(a, b) \neq (0, 0)$, is $R > 0$. Het punt $\bigl(\frac{a}{R}, \frac{b}{R}\bigr)$ ligt op de eenheidscirkel, want
\bewijsregel{\Bigl(\frac{a}{R}\Bigr)^2 + \Bigl(\frac{b}{R}\Bigr)^2 = \frac{a^2 + b^2}{R^2} = 1 .}
Er is dus een hoek $\varphi$ met $\cos\varphi = \frac{a}{R}$ en $\sin\varphi = \frac{b}{R}$. Met de somformule voor de sinus is dan
\bewijsregel{R\sin(x + \varphi) = R\cos\varphi\,\sin x + R\sin\varphi\,\cos x = a\sin x + b\cos x .}
\end{bewijs}
\end{theorie}

De som $a\sin x + b\cos x$ schommelt dus tussen $-R$ en $R$. Haar grootste waarde is $R = \sqrt{a^2 + b^2}$, en die is meestal kleiner dan $|a| + |b|$: de toppen van de twee golven vallen niet samen.

\voorbeeld{}{Schrijf $\sin x + \sqrt3\cos x$ als $R\sin(x + \varphi)$. Wat is de grootste waarde?\label{vb:h2-samenvoegen}}{%
\uitleg{Hier is $a = 1$ en $b = \sqrt3$. We berekenen $R$, en zoeken dan de hoek met cosinus $\frac{a}{R}$ en sinus $\frac{b}{R}$.}
\begin{stappenbreed}
R = \sqrt{1^2 + (\sqrt3)^2} = \sqrt{4} = 2 & \\[1ex]
\cos\varphi = \tfrac12,\ \sin\varphi = \tfrac{\sqrt3}{2} \iff \varphi = \tfrac{\pi}{3} & Beide positief: eerste kwadrant. Bijzondere waarden. \\[1ex]
\sin x + \sqrt3\cos x = 2\sin\bigl(x + \tfrac{\pi}{3}\bigr) & De grootste waarde is $2$.
\end{stappenbreed}}

Is $\varphi$ geen hoek uit de tabel, dan zoek je hem met het rekentoestel. Pas dan op: de toets $\sin^{-1}$ geeft altijd een hoek tussen $-\frac{\pi}{2}$ en $\frac{\pi}{2}$, met een positieve cosinus. Dat is niet altijd de juiste.

\voorbeeld{}{Schrijf $-3\sin x + 4\cos x$ als $R\sin(x + \varphi)$, met $\varphi$ op vier decimalen.\label{vb:h2-samenvoegen2}}{%
\uitleg{Nu is $a = -3$ en $b = 4$. Uit de tekens van $\cos\varphi$ en $\sin\varphi$ halen we eerst het kwadrant. Pas dan zoeken we de hoek.}
\begin{stappenbreed}
R = \sqrt{9 + 16} = 5 & \\[1ex]
\cos\varphi = -\tfrac35,\ \sin\varphi = \tfrac45 & Cosinus negatief, sinus positief: $\varphi$ ligt in het tweede kwadrant. \\[1ex]
\sin^{-1}(0{,}8) \approx 0{,}9273 & Het rekentoestel geeft een hoek in het eerste kwadrant, met dezelfde sinus. \\[1ex]
\varphi = \pi - 0{,}9273 \approx 2{,}2143 & De supplementaire hoek heeft dezelfde sinus, en ligt wel in het tweede kwadrant. \\[1ex]
-3\sin x + 4\cos x \approx 5\sin(x + 2{,}2143) &
\end{stappenbreed}
Controle met $x = 0$: links staat $4$, rechts $5\sin 2{,}2143 \approx 4{,}000$.}

\begin{werkwijze}
Bereken $R = \sqrt{a^2 + b^2}$. Bepaal het kwadrant van $\varphi$ uit de tekens van $a$ en $b$. Zoek $\varphi$ met de tabel, of met $\sin^{-1}\frac{b}{R}$ op het rekentoestel, en neem zo nodig de supplementaire hoek. Controleer door $x = 0$ in te vullen.
\end{werkwijze}

\begin{oefeningen}
\oef{h2-samenvoegen}

\oef{h2-uitwerkk}
\end{oefeningen}

%========================================================================================
\opdrachtenkop

\begin{oefeningen}
\oef{h2-o1}

\oef{h2-o13}

\oef{h2-o10}

\oef{h2-o2}

\oef{h2-o14}

\oef{h2-o12}

\oef{h2-o11}

\oef{h2-o3}

\oef{h2-o15}

\oef{h2-o16}

\oef{h2-o18}

\oef{h2-o4}

\oef{h2-o23}

\oef{h2-o17}

\oef{h2-o19}

\oef{h2-o5}

\oef{h2-o21}

\oef{h2-o22}

\oef{h2-o6}

\oef{h2-o20}

\oef{h2-o7}

\oef{h2-o8}

\oef{h2-o24}

\oef{h2-o27}

\oef{h2-o9}

\oef{h2-o26}

\oef{h2-o25}
\end{oefeningen}

%----------------------------------------------------------------------------------------
% KOPPEL: verbanden met andere leerstof
%----------------------------------------------------------------------------------------
\needspace{20\baselineskip}
\opdrachtsectie{Koppel}

\lkr{Integralen en Matrices: kijk na wat al gezien is. Het pakket legt de volgorde niet vast.}

\begin{oefeningen}

\oef{h2-koppel-2}

\oef{h2-koppel-3}
\end{oefeningen}

%----------------------------------------------------------------------------------------
\opdrachtsectie{Test jezelf}

\begin{oefeningen}
\oef{h2-t1}

\oef{h2-t2}

\oef{h2-t3}

\oef{h2-t6}

\oef{h2-t4}

\oef{h2-t7}

\oef{h2-t5}
\end{oefeningen}

\herhaalschema{%
  \fv{oef:h2-t1}{vb:h2-cossom}{oef:h2-cossom-k}%
  \fv{oef:h2-t2}{vb:h2-som}{oef:h2-cosverschil-k}%
  \fv{oef:h2-t3}{vb:h2-dubbel}{oef:h2-verdubbel}%
  \fv{oef:h2-t6}{vb:h2-carnot}{oef:h2-carnotk}%
  \fv{oef:h2-t4}{vb:h2-identiteit}{oef:h2-identk}%
  \fv{oef:h2-t7}{vb:h2-tansom}{oef:h2-o6}%
  \fv{oef:h2-t5}{vb:h2-samenvoegen}{oef:h2-samenvoegen}}

\ifzelfstudie
\section*{Maak je eigen samenvatting}
Sluit dit boek en schrijf in je schrift: de som- en verschilformules (met het idee van het bewijs), de verdubbelingsformules met de drie vormen van $\cos 2\alpha$, de formules van Carnot, en hoe je $a\sin x + b\cos x$ samenvoegt. Vergelijk daarna met de kaders.
\fi
